\section{Orchestration and collaboration}

% ── slides p.2–p.4 ──────────────────────────────────────────
\subsection{Why diagrams?}

Graphical languages are first of all a \emph{communication} medium. Their value comes from a careful selection of symbols (shapes, colours, arrows) assembled into an alphabet shared by the stakeholders: intuitive enough for a domain expert without mathematical training, precise enough to be unambiguous. Expressivity does not justify clutter, though: an auto-generated flowchart that wires dozens of decorative boxes (\textsc{start}, \textsc{decision}, \textsc{database}, \textsc{document}, \dots) into a baroque graph is unreadable. A good process diagram uses \emph{few} graphical elements and a small number of clear connections.

% ── slide p.5 ──────────────────────────────────────────
The three notations seen in the course share the ingredient catalogue but disagree on the symbols:
\begin{center}
\begin{tabular}{lccc}
\toprule
 & \textbf{BPMN} & \textbf{EPC} & \textbf{WfN} \\
\midrule
start / end & circle      & hexagon      & circle (place)        \\
join / split & diamond     & circle      & square (transition)    \\
tasks        & rounded rect.& rounded rect.& square (transition)   \\
links        & solid arrow & dashed arrow & solid arrow           \\
\bottomrule
\end{tabular}
\end{center}

% ── slides p.6–p.7 ──────────────────────────────────────────
The insurance-claim process of the previous chapters, redrawn in EPC syntax with tools such as ARIS Express\footnote{\url{https://ariscommunity.com/aris-express}} or VP Online,\footnote{\url{https://online.visual-paradigm.com/}} keeps the same semantic content as its BPMN rendering: activities (\textsc{recording}, \textsc{type}, \textsc{policy}, \textsc{premium}, \dots, \textsc{filing}) alternate with events, and small AND/XOR connector circles realise the splits and joins. EPC layouts run conventionally top to bottom and come out taller than their BPMN counterparts.

% ── slides p.8–p.10 ──────────────────────────────────────────
\subsection{Event-driven process chains (EPC)}

% Corrected name order: the ARIS author is August-Wilhelm Scheer.
The EPC method was developed in the early 1990s by August-Wilhelm Scheer as the control-flow component of the \textbf{ARIS framework} (Architecture of Integrated Information Systems), a holistic approach to enterprise modelling covering organisation, data, function, and control views.\footnote{Scheer A.-W., \emph{ARIS: Vom Gesch\"aftsprozess zum Anwendungssystem}, Springer, 1998; \emph{ARIS: Business Process Modeling}, 2nd ed., Springer, 2000.} EPC is used to configure ERP implementations, to drive the modelling of business processes, and to support their analysis and redesign. The notation is informal (no legend is needed to read a diagram) and serialises into the XML-based \textbf{EPC Markup Language} (\texttt{.epml}).

% ── slide p.11 ──────────────────────────────────────────
\begin{definitionbox}{Event-driven Process Chain (EPC)}
An \textbf{Event-driven Process Chain} (EPC) is an ordered graph of \textbf{events} and \textbf{functions}. It provides various \textbf{connectors} that allow alternative and parallel execution of processes; the connectors are further qualified by the usage of \textbf{logical operators} such as OR, AND, and XOR.
\end{definitionbox}

The claimed pedagogical strength of EPC is simplicity: regardless of formal expressiveness, the easy-to-understand look of the diagrams made it a widely accepted industrial notation.

% ── slides p.12–p.18 ──────────────────────────────────────────
\subsubsection{Ingredients}

\begin{center}
\begin{tikzpicture}[node distance=1.0cm]
  \node[draw=accent, thick, regular polygon, regular polygon sides=6, minimum width=1.5cm, font=\scriptsize] (ev) {event};
  \node[draw=accent, thick, rounded corners=4pt, fill=rowgray, minimum width=1.6cm, minimum height=0.9cm, font=\scriptsize, right=of ev] (fn) {function};
  \node[draw=accent, thick, circle, minimum size=0.9cm, font=\scriptsize, right=of fn] (and) {$\wedge$};
  \node[draw=accent, thick, circle, minimum size=0.9cm, font=\scriptsize, right=2pt of and] (xor) {XOR};
  \node[draw=accent, thick, circle, minimum size=0.9cm, font=\scriptsize, right=2pt of xor] (or) {$\vee$};
  \draw[->, dashed, thick, accent] ($(or.east)+(0.5,0)$) -- node[above, font=\scriptsize] {flow} ++(1.6,0);
\end{tikzpicture}
\end{center}

\begin{itemize}
  \item \textbf{Events} (hexagons) are the passive elements: they describe the circumstances under which a function works or the state it results in: pre- and post-conditions in a Hoare-style reading. Every EPC must \emph{start and end with events}, and along the flow events and functions \emph{alternate}: between any two functions sits at least one event, and vice versa.
  \item \textbf{Functions} (rounded rectangles) are the active elements: the tasks of the business process. A function can be refined into another EPC, supporting hierarchical modelling as in workflow nets.
  \item \textbf{Connectors} (circles, octagons in some tools) carry one of three logical operators, often abbreviated by the propositional symbols $\wedge = \mathrm{AND}$, $\times = \mathrm{XOR}$, $\vee = \mathrm{OR}$: AND synchronises parallel branches, XOR routes an exclusive choice, OR expresses inclusive choice (one or more branches, possibly all).
  \item \textbf{Control-flow connections} (dashed arrows) link events to functions or connectors and vice versa, expressing causal dependency: the source must occur before the target.
\end{itemize}

% ── slide p.19 ──────────────────────────────────────────
A connector with one incoming arc and several outgoing arcs is a \textbf{split}; with several incoming arcs and one outgoing arc it is a \textbf{join}. Combined with the three operators this gives six patterns (AND/XOR/OR-split, AND/XOR/OR-join); the arrow geometry alone disambiguates split from join, so the glyph does not change.\footnote{Weske M., \emph{Business Process Management}, Springer, 2007.}

% ── slides p.20–p.21 ──────────────────────────────────────────
\begin{notebox}{OR-join: intended meaning}
The OR-join is the most subtle connector:
\begin{itemize}
  \item if only one input arrives, it releases the flow;
  \item if both inputs arrive, it releases exactly one output (not two);
  \item if one input arrives, it must wait until the other arrives \emph{or} until it is guaranteed that the other will never arrive.
\end{itemize}
The last clause is what makes OR-joins formally hard: deciding that ``the other will never arrive'' requires reasoning about the global state, which a local firing rule cannot do. This is the source of the well-known semantic difficulties around OR-joins.
\end{notebox}

The OR-join is not a luxury. When two upstream branches may both produce a token, both omit it, or only one of them produces one, the OR-join is the only construct that merges them correctly: an XOR-join would expect exactly one input and break on the simultaneous arrival of both, an AND-join would wait forever when one branch stays silent.

% ── slide p.22 ──────────────────────────────────────────
\subsubsection{A taste of \texttt{.epml}}

EPML is the XML interchange format for EPC.\footnote{Mendling J., N\"uttgens M., ``EPC markup language (EPML): an XML-based interchange format for event-driven process chains (EPC)'', \emph{Information Systems and e-Business Management}, 4(3):245--263, Springer, 2006.} A file is rooted at an \verb|<epml:epml>| element; a \verb|<directory>| holds \verb|<epc>| containers carrying the events, functions, connectors, and arcs. An extract of an ``Online Shopping'' EPC:
\begin{verbatim}
<?xml version="1.0" encoding="UTF-8"?>
<epml:epml xmlns:epml="http://www.epml.de">
  <directory name="Root">
    <epc epcId="1" name="Online Shopping">
      <event id="1"><name>Start Online Shopping</name></event>
      <arc id="10"><flow source="1" target="2"/></arc>
      <function id="2"><name>Determine List of Products</name></function>
      <arc id="11"><flow source="2" target="3"/></arc>
      <event id="3"><name>Not all Products in Cart</name></event>
      <arc id="12"><flow source="3" target="4"/></arc>
      <xor id="4"/>
      ...
      <processInterface id="111">
        <name>Continue with Order Process</name>
      </processInterface>
    </epc>
  </directory>
</epml:epml>
\end{verbatim}
The event--arc--function alternation is preserved in the XML, and a \verb|<processInterface>| pointer lets a process call another process: the EPML counterpart of a sub-process.\footnote{\url{http://www.epml.de}}

% ── slides p.23–p.25 ──────────────────────────────────────────
\subsection{Orchestration}

The rest of the chapter develops a buyer--reseller pair taken from Weske.\footnote{Weske M., \emph{Business Process Management}, Springer, 2007, Sect.~1.1.}

\begin{definitionbox}{Orchestration}
Business process models are by definition enacted within a \emph{single} organisation. The ordering of activities can therefore be controlled by a \textbf{business process management system} run \emph{centrally} by that organisation. This kind of control is called \textbf{orchestration}.
\end{definitionbox}

The name comes from the conductor of an orchestra. An orchestration is a \emph{single-viewpoint} model: one organisation owns the process and dictates the flow.

% ── slides p.26–p.29, p.32–p.33 ──────────────────────────────────────────
The reseller's orchestration contains \textsc{Receive Order}, followed by an AND-split into a billing branch (\textsc{Send Invoice} $\to$ \textsc{Receive Payment}) and a shipping branch (\textsc{Ship Products}), synchronised before completion (in the EPC rendering, before an \textsc{Archive Order} function). The same single-viewpoint process can be drawn in all three notations (EPC with alternating events, as a workflow net whose places hold the case state, or as a BPMN pool) and the reseller uses it to configure its own BPM system: every instance then runs as specified, with billing and shipping concurrent. Symmetrically, the buyer owns an orchestration: \textsc{Place Order}, then concurrently \textsc{Receive Products} and \textsc{Receive Invoice} $\to$ \textsc{Settle Invoice}, synchronised before completion. In both diagrams no message to or from the other participant appears.

% ── slides p.30–p.31, p.34–p.36 ──────────────────────────────────────────
\subsection{Collaboration}

\begin{definitionbox}{Collaboration}
Each business process is enacted by one organisation, but cross-organisational business processes can \emph{interact} with each other. Linking the activities of single business processes that are related is called \textbf{collaboration}. A collaboration describes a \emph{multi-point}, \emph{decentralised} viewpoint: it comprises the activities of autonomous business processes and how they interact to achieve a common goal.
\end{definitionbox}

There is no longer a single conductor: separately developed processes must communicate, exchanging information of any kind: electronic messages, physically transported objects (the shipped products), documents (the invoice), and the model must make the exchanges explicit so that the global behaviour can be reasoned about. In BPMN the pools sit side by side and the exchanges become \textbf{message flows}, drawn as dashed arcs. Four of them weave the two orchestrations into one collaboration:
\begin{center}
\resizebox{0.92\linewidth}{!}{%
\begin{tikzpicture}
  \node[bpmn event] (bs) at (0.6,3.2) {};
  \node[bpmn task] (po) at (2.0,3.2) {Place Order};
  \node[bpmn gateway] (bg1) at (3.7,3.2) {$+$};
  \node[bpmn task] (ri) at (5.6,3.9) {Receive Invoice};
  \node[bpmn task] (si) at (8.2,3.9) {Settle Invoice};
  \node[bpmn task] (rp) at (6.9,2.5) {Receive Products};
  \node[bpmn gateway] (bg2) at (10.1,3.2) {$+$};
  \node[bpmn event, thick] (be) at (11.3,3.2) {};
  \draw[bpmn flow] (bs) -- (po); \draw[bpmn flow] (po) -- (bg1);
  \draw[bpmn flow] (bg1) |- (ri); \draw[bpmn flow] (ri) -- (si);
  \draw[bpmn flow] (bg1) |- (rp);
  \draw[bpmn flow] (si) -| (bg2); \draw[bpmn flow] (rp) -| (bg2);
  \draw[bpmn flow] (bg2) -- (be);
  \draw[accent] (-0.2,1.8) rectangle (12.1,4.6);
  \node[font=\small\bfseries, text=accent, anchor=north west] at (-0.2,4.6) {Buyer};
  \node[bpmn event] (rs) at (0.6,-0.6) {};
  \node[bpmn task] (ro) at (2.0,-0.6) {Receive Order};
  \node[bpmn gateway] (rg1) at (3.7,-0.6) {$+$};
  \node[bpmn task] (sv) at (5.6,0.1) {Send Invoice};
  \node[bpmn task] (rpay) at (8.2,0.1) {Receive Payment};
  \node[bpmn task] (sp) at (6.9,-1.3) {Ship Products};
  \node[bpmn gateway] (rg2) at (10.1,-0.6) {$+$};
  \node[bpmn event, thick] (re) at (11.3,-0.6) {};
  \draw[bpmn flow] (rs) -- (ro); \draw[bpmn flow] (ro) -- (rg1);
  \draw[bpmn flow] (rg1) |- (sv); \draw[bpmn flow] (sv) -- (rpay);
  \draw[bpmn flow] (rg1) |- (sp);
  \draw[bpmn flow] (rpay) -| (rg2); \draw[bpmn flow] (sp) -| (rg2);
  \draw[bpmn flow] (rg2) -- (re);
  \draw[accent] (-0.2,-2.0) rectangle (12.1,0.8);
  \node[font=\small\bfseries, text=accent, anchor=south west] at (-0.2,-2.0) {Reseller};
  \draw[bpmn flow, dashed] (po.south) -- (ro.north);
  \draw[bpmn flow, dashed] (sv.north) -- (ri.south);
  \draw[bpmn flow, dashed] (sp.north) -- (rp.south);
  \draw[bpmn flow, dashed] (si.south) -- (rpay.north);
\end{tikzpicture}}
\end{center}
The order, the invoice, the products, and the payment each cross the pool boundary exactly once; what one pool sends must match what the other expects.

% ── slides p.37–p.40 ──────────────────────────────────────────
\subsection{Choreography}

\begin{definitionbox}{Choreography}
The interactions of a set of business processes can be specified also as a process \textbf{choreography}: a \emph{global-viewpoint} model that contains only the activities related to interactions between participants.
\end{definitionbox}

The choreography is distinguished by what it removes: unlike an orchestration it has \emph{no central agent} controlling the activities, and unlike a collaboration it contains \emph{only} the interaction activities, not the internal ones. The name comes from dancers: each behaves autonomously but follows the prescribed steps and timing. For the buyer--reseller pair the choreography lists only four interactions: \textsc{Buyer places the order to Reseller}, \textsc{Reseller ships products to Buyer}, \textsc{Reseller sends the invoice to Buyer}, \textsc{Buyer settles the invoice to Reseller}. A choreography is a sort of \emph{contract}: it states how actors \emph{should} interact, not how each realises its part. The three viewpoints do not subsume one another:
\begin{itemize}
  \item \textbf{orchestration}, single viewpoint: one organisation, one pool (for the implementation team);
  \item \textbf{collaboration}, multiple viewpoints: several pools, each running its own orchestration, wired by message flows (for the integration team);
  \item \textbf{choreography}, global viewpoint: only the message exchanges, read as a contract (for inter-organisation negotiation).
\end{itemize}

% ── slides p.41–p.50 ──────────────────────────────────────────
\subsection{Compatibility}

Different orchestrations of the same participant are possible, but do they all make sense, and do they interact correctly with the counterpart? Two reseller and four buyer variants set up the question.\footnote{Weske M., \emph{Business Process Management}, Springer, 2007.}
\begin{itemize}
  \item \textbf{$R_1$}: \textsc{Receive Order} $\to$ AND-split into \{\textsc{Send Invoice} $\to$ \textsc{Receive Payment}\} and \{\textsc{Ship Products}\}.
  \item \textbf{$R_2$}: \textsc{Receive Order} $\to$ \textsc{Send Invoice} $\to$ \textsc{Receive Payment} $\to$ \textsc{Ship Products} (fully sequential: payment before shipping).
  \item \textbf{$B_1$}: \textsc{Place Order} $\to$ AND-split into \{\textsc{Receive Invoice} $\to$ \textsc{Settle Invoice}\} and \{\textsc{Receive Products}\}.
  \item \textbf{$B_2$}: \textsc{Place Order} $\to$ AND-split into \{\textsc{Receive Products} $\to$ \textsc{Settle Invoice}\} and \{\textsc{Receive Invoice}\}: pays only after the goods arrive.
  \item \textbf{$B_3$}: \textsc{Place Order} $\to$ AND-split into \{\textsc{Receive Products}\} and \{\textsc{Receive Invoice}\}, AND-join, then \textsc{Settle Invoice}: pays only after \emph{both} goods and invoice.
  \item \textbf{$B_4$}: \textsc{Place Order} $\to$ \textsc{Receive Invoice} $\to$ \textsc{Receive Products} $\to$ \textsc{Settle Invoice} (strict sequence).
\end{itemize}

% ── slides p.51–p.58 ──────────────────────────────────────────
\begin{examplebox}{Exercise: compatibility matrix}
Build the $2 \times 4$ matrix marking every pair $(R_i, B_j)$ whose interaction may run into problems: an activity fired in an order the counterpart is not ready to accept. Optionally extend the matrix with further variants.
\end{examplebox}

Pairwise reasoning fills the matrix:
\begin{itemize}
  \item $R_1 + B_1$: \textbf{ok}: both sides handle billing and shipping concurrently; every message finds its receiver.
  \item $R_1 + B_2$: \textbf{ok}: $R_1$ ships without waiting for payment, so $B_2$ gets the products, settles, and $R_1$ collects the payment.
  \item $R_1 + B_3$: \textbf{ok}: $R_1$ sends invoice and products concurrently; $B_3$ receives both, then settles.
  \item $R_1 + B_4$: \textbf{ok}: the AND-split of $R_1$ is permissive enough to serve the buyer's strict receive-invoice-then-products sequence.
  \item $R_2 + B_1$: \textbf{ok}: $B_1$ settles upon receiving the invoice without waiting for the products, so $R_2$ gets its payment and ships last.
  \item $R_2 + B_2$: \textbf{no}: $R_2$ ships only after the payment, $B_2$ pays only after the products: each side waits for the other and the collaboration deadlocks.
  \item $R_2 + B_3$: \textbf{no}: same deadlock: $B_3$ needs the products before settling, $R_2$ needs the payment before shipping.
  \item $R_2 + B_4$: \textbf{no}: $B_4$ wants the products before settling; $R_2$ again waits for the payment first.
\end{itemize}
\begin{center}
\begin{tabular}{c|cccc}
      & $B_1$ & $B_2$ & $B_3$ & $B_4$ \\
\hline
$R_1$ & ok    & ok    & ok    & ok    \\
$R_2$ & ok    & no    & no    & no    \\
\end{tabular}
\end{center}
The pattern is informative. $R_1$, the AND-split reseller, is \emph{permissive} and composes with all four buyers; $R_2$, the strict-sequence reseller, only works with buyers willing to pay before the goods arrive. Two locally-correct orchestrations can compose into a deadlocking collaboration: compatibility is a genuinely global property, formalised later in the course.

% ── slides p.59–p.64 ──────────────────────────────────────────
\subsection{Exercises: the travel agency}

\begin{examplebox}{Travel-agency orchestration (single viewpoint)}
Design a travel-agency orchestration in BPMN, EPC, and WfN. The tasks handle a travel request and include booking a flight, booking a hotel, changing dates, cancelling the booking, and confirming the booking. Draw a process diagram relating the tasks.
\end{examplebox}

% Corrected: the initial XOR is the loop join, not a Book-Car choice (no such
% activity exists in the exercise or slide p.60).
\textit{Solution sketch.}

In EPC: start event \textsc{Travel request} flows into an XOR connector that also collects the \textsc{Change dates} return branch and passes to the function \textsc{Receive dates}. An AND connector then forks \textsc{Book Hotel} and \textsc{Book Flight} and joins them again. A downstream XOR lets the customer \textsc{Change dates} (loop back to the first XOR), \textsc{Confirm} (end event \textsc{Success}), or \textsc{Cancel} (end event \textsc{Failure}).

The BPMN version replaces connectors with gateways (XOR diamonds for the decisions, AND diamond for the fork) inside a single pool with two end events.

The WfN version realises the same structure: competing transitions on a shared input place for each XOR, one transition with two output places for the AND, the change branch looping back, and the sink reached through either the confirm or the cancel transition. More verbose, but unambiguous.

% ── slides p.65–p.66 ──────────────────────────────────────────
\begin{examplebox}{Travel choreography (global viewpoint)}
Model the same scenario (booking flight and hotel with the possibility to change dates, cancel, or confirm) as a BPMN choreography between \textsc{Tourist} and \textsc{Travel agency}.
\end{examplebox}

\textit{Solution sketch.} Only the message exchanges appear: \textsc{Travel request} (tourist $\to$ agency), \textsc{Travel planned} (agency $\to$ tourist), \textsc{Ask confirmation} (agency $\to$ tourist), then the choice between \textsc{Change dates} (loop back), \textsc{Confirm} (tourist $\to$ agency, success), and \textsc{Cancel} (failure). No internal activity of either side is shown.

% ── slides p.67–p.71 ──────────────────────────────────────────
\begin{examplebox}{Tourist orchestration and full collaboration}
Draw the tourist's own orchestration (BPMN, EPC, WfN), then assemble the BPMN collaboration of tourist and travel agency.
\end{examplebox}

\textit{Solution sketch.} The tourist's pool: \textsc{Issue travel request}, a message-receive task for the planned trip, then an XOR decision among \textsc{Confirm}, \textsc{Cancel}, and \textsc{Change dates} (looping back to wait for a new plan), ending in a success or failure event. The collaboration places the two pools side by side, each running its orchestration, with one message flow per choreography interaction: \textsc{Travel request} and \textsc{Confirm}/\textsc{Cancel} flow tourist $\to$ agency, \textsc{Travel planned} and \textsc{Ask confirmation} flow agency $\to$ tourist, and the \textsc{Change dates} message returns the agency to its booking phase, restarting the negotiation.
